\documentclass[11pt]{article}

\usepackage[a4paper,margin=1in]{geometry}
\usepackage{amsmath,amssymb}
\usepackage{booktabs}
\usepackage{xcolor}
\usepackage{enumitem}
\usepackage[colorlinks=true,linkcolor=blue!50!black,urlcolor=blue!50!black]{hyperref}
\usepackage{titlesec}
\titlespacing{\section}{0pt}{1.2em}{0.5em}
\titlespacing{\subsection}{0pt}{0.9em}{0.3em}

\newcommand{\rsix}{\texttt{responder\_6}}
\newcommand{\Rw}{R^2_w}

\title{\vspace{-1.5em}Jane Street Real-Time Market Data Forecasting\\[0.3em]
\large Models and Angles: A Technical Report}
\author{Alberto Gennaro}
\date{\today}

\begin{document}
\maketitle
\vspace{-1.5em}

\begin{abstract}
\noindent
We document the modelling approaches taken on the Jane Street Real-Time Market
Data Forecasting task: predicting the weighted-$R^2$ of \rsix{} from 79
anonymized features streamed per \texttt{(symbol, date, time)}. We first
summarise what the target \emph{is} --- a forward-shifted moving average of a
near-white-noise signal, which caps achievable skill --- and then describe the
leakage-safe data pipeline, the twelve models we tried across five
architectural families, the signature and covariate ideas, and the ensembling
strategy. The best result to date is a three-model bagged average scoring
$\Rw = +0.01039$ on a held-out 100-day tail, versus the $+0.0112$ that took
8th place. The central empirical lesson is that, because the target is
dominated by noise, \emph{variance reduction} (online refitting, temporal
bagging, ensembling) buys more than any single stronger model.
\end{abstract}

\section{Problem and the nature of the target}

Each row is indexed by \texttt{(symbol\_id, date\_id, time\_id)} with 79
features and 9 responders; we predict \rsix{}, scored by the weighted, zero-mean
$R^2$
\[
\Rw = 1 - \frac{\sum_i w_i (\hat y_i - y_i)^2}{\sum_i w_i\, y_i^2}.
\]

\paragraph{What the responders are.}
Autocorrelation and Fourier analysis (Kaggle discussion \#555562, corroborated
by our own tag analysis) show the responders are \emph{moving averages of a few
underlying signals}, shifted forward in time:
\rsix{} is an SMA of window 20, \texttt{responder\_7} of window 120,
\texttt{responder\_8} of window 4, of one signal (``exchange A''); responders
$\{3,4,5\}$ are the same windows of a second signal (``exchange B''); and
$\{0,1,2\}$ are their differences. The label at time $t$ is effectively the
SMA-20 over $[t, t{+}20]$. The practical consequence is stark: for a forward
moving average of near-white noise, $\mathbb{E}[\text{future}\mid\text{present}]
\approx 0$, so the predictable component is tiny. This explains why all scores,
ours included, plateau near $+0.010$, and it motivates our emphasis on variance
reduction over single-model capacity.

\paragraph{Deployability constraint.}
The evaluation gateway streams one \texttt{time\_id} at a time and exposes
responders only as \emph{daily lags} (the previous day's full path). Hence
(i)~any feature used at time $t$ may look back only to $\le t$ within the day,
and (ii)~the target's own history is available only as yesterday's lagged
responders. Every model we build respects this; we verify it explicitly.

\section{Data pipeline and leakage control}

\begin{itemize}[leftmargin=1.2em,itemsep=1pt]
  \item \textbf{Date range.} The per-day time grid stabilises at 968 steps from
  \texttt{date\_id}~677; following the 8th-place solution we train from date~700.
  Training pool: 700--1598; held-out validation tail: 1599--1698.
  \item \textbf{Features.} Raw features; per-symbol trailing rolling mean/std
  (window 1000); cross-sectional market average per \texttt{(date,time)};
  two synthetic auxiliary responders; optionally previous-day lagged responders
  and signals derived from them.
  \item \textbf{Leakage audit.} Model inputs are trailing or contemporaneous
  (no forward look). The auxiliary targets use forward shifts by construction;
  their only cross-boundary reach is removed by a one-day \emph{embargo} between
  the pool and the validation tail. Standardization statistics are fit on
  training dates only; the online scaler adapts on past days only. Lagged
  responders are strictly previous-day (a \texttt{date+1} self-join), verified
  backward-only.
  \item \textbf{Engineering for scale.} Features are pre-computed once for the
  whole pool into a \texttt{Float16} memory-mapped array; each model then reads
  only the dates it trains on. This is what makes on-laptop training of many
  models on hundreds of dates feasible.
\end{itemize}

\section{Models and the angle behind each}

We treat one \texttt{(symbol, day)} as one training example: a sequence of 968
timesteps. Predictions are per-timestep and strictly causal within the day.

\subsection{Gradient boosting (\texttt{xgb})}
A per-row \texttt{XGBRegressor} on the flat feature matrix --- no time
structure. Cheap, strong, and a natural ensemble base because its errors are
uncorrelated with the sequence models. Best single static score $\approx
+0.0068$.

\subsection{Feed-forward baselines (\texttt{mlp})}
A plain MLP on the per-row features, as a sanity floor. It scores
\emph{negative} ($-0.0026$): the signal here is temporal, and a model with no
access to the path cannot recover it. Useful precisely as a negative control.

\subsection{Recurrent models and online refitting (\texttt{gru}, \texttt{lstm}, \texttt{*\_modelr})}
GRU/LSTM encoders over the daily sequence. The \texttt{modelr} variants
replicate the 8th-place design: four parallel recurrent branches plus auxiliary
heads that also predict responders 7--10, a multi-task regulariser. The
decisive ingredient is \emph{online refitting}: at inference the model takes one
gradient step per day on the previous day's realised responders (using the
daily lags), adapting to drift. This lifts RNN scores by $+0.003$ to $+0.005$
and is the single most valuable trick we found. \texttt{lstm\_modelr} with
refit learning rate $10^{-3}$ is our best single model at $\Rw = +0.0092$.

\subsection{Volterra signatures (\texttt{sig\_transformer}, \texttt{mlp\_sig})}
Motivated by rough-path theory, we augment models with the truncated
\emph{Volterra signature} of a small set of channels over a rolling within-day
window: increments are weighted by a power-law kernel $(T-s)^{H-1/2}$ (Hurst
$H<1/2$ emphasises recent moves), and the truncated tensor algebra is computed
by Chen iteration. We explored (i)~classical vs \emph{log}-signatures --- the
latter live in the free Lie algebra and are lower-dimensional and less
redundant; the minimal Lyndon-basis log-signature (via \texttt{iisignature})
cuts depth-3 dimension from 819 to 285; and (ii)~\emph{channel selection}: an
AR($p$) ridge fit scoring how well each feature's own lags predict the target,
which cleanly beat a naive univariate-correlation filter. A stateless
per-timestep MLP fed the minimal log-signature reached $+0.0066$ online ---
competitive with the transformer while far simpler --- but signature streams
never improved the ensemble, being too correlated with the RNNs.

\subsection{Transformers: vanilla, inverted, and cross-attention (\texttt{transformer}, \texttt{itransformer}, \texttt{timexer})}
A causal attention encoder is the baseline. We then adapted two recent ideas:
\begin{itemize}[leftmargin=1.2em,itemsep=1pt]
  \item \textbf{iTransformer} (inverted attention): attention runs \emph{across
  features} rather than time, modelling cross-feature correlations directly. We
  make it causal by applying the inverted attention per timestep and mixing time
  with a separate causal block.
  \item \textbf{TimeXer} (exogenous covariates): separate an \emph{endogenous}
  target series from \emph{exogenous} covariates, with a learnable global token
  bridging them by cross-attention. In our mapping the endogenous stream is
  \emph{yesterday's responder path} (fully known at day start, hence patchable
  non-causally) and the exogenous stream is today's features (causal). Because
  turning today's features into whole-window variate tokens would leak the
  future, we invert the bridge direction relative to the paper --- the causal
  today-stream queries the fully-known yesterday-stream. Causality is verified
  by a perturbation test.
\end{itemize}

\section{Feature-engineering angles}
Beyond the baseline set we investigated: (i)~\textbf{lagged responders} ---
yesterday's full responder path as input features, which is both what the
gateway provides and the natural ``endogenous'' signal; (ii)~\textbf{responder
signals} derived from them --- multi-scale momentum (differences across the
SMA-4/20/120 windows) and the cross-exchange spread, motivated directly by the
reverse-engineered structure; and (iii)~\textbf{tag-group structure} from
\texttt{features.csv}, where matched feature triples differing only in tags
\{12,13,14\} behave like three views of one quantity, suggesting level/slope
aggregates. Each is added as an opt-in group and validated by ablation before
adoption.

\section{Ensembling and temporal bagging}
Given the noise ceiling, combining models matters more than any single one. A
plain average of \texttt{xgb} $+$ \texttt{gru\_modelr}(online) $+$
\texttt{lstm\_modelr}(online) scores $\Rw = +0.01039$, about 93\% of the
8th-place mark, and beats both a convex weight-optimised blend and a non-linear
XGB stacker (which overfit the fitting split). Adding signature or transformer
streams did \emph{not} help --- they are too correlated with the RNNs. The
response is \emph{temporal bagging}: train the same architecture on different
overlapping date resamples (subsampling $\sim$60\% of the pool per member,
different seeds), producing decorrelated members whose average a simple mean can
exploit. This is now the primary lever.

\section{Results}

\begin{table}[h]
\centering
\small
\begin{tabular}{lrr}
\toprule
Model (280 train dates unless noted) & static $\Rw$ & online $\Rw$ \\
\midrule
\texttt{lstm\_modelr} (refit lr $10^{-3}$)      & $+0.00637$ & $\mathbf{+0.00920}$ \\
\texttt{gru\_modelr}                            & $+0.00763$ & $+0.00893$ \\
\texttt{xgb}                                    & $+0.00684$ & --- \\
\texttt{mlp\_sig} (minimal log-sig, AR channels)& $+0.00614$ & $+0.00657$ \\
\texttt{sig\_transformer} (AR ch., refit $3\!\times\!10^{-5}$) & $+0.00594$ & $+0.00623$ \\
\texttt{mlp} (non-sequence control)             & $-0.00258$ & --- \\
\midrule
\textbf{3-model simple-average ensemble}        & \multicolumn{2}{c}{$\mathbf{+0.01039}$} \\
\bottomrule
\end{tabular}
\caption{Weighted-$R^2$ on the 1599--1698 validation tail. 8th-place private LB
was $+0.0112$.}
\end{table}

\paragraph{Selected findings.}
(1)~More training data helped every model, whereas doubling model width did not
--- the bottleneck is data and noise, not capacity. (2)~Online refit is the RNN
lever; transformers barely benefit and signature-transformers need a
$30\times$ smaller refit LR. (3)~A static-XGB ablation (200 train days)
found that adding the nine lagged responders as \emph{flat} features did not
help ($+0.00594 \to +0.00591$), but the five \emph{engineered} responder
signals --- multi-scale momentum and the venue spread, differences across the
SMA-4/20/120 windows --- added $+0.00030$ ($+0.00594 \to +0.00624$, a $\sim$5\%
relative gain). The interpretation is instructive: a tree cannot cheaply
synthesise cross-window differences from the raw lags, so supplying them
explicitly pays off, whereas the raw lags on their own are near-useless to a
static model --- consistent with the noise ceiling. The lags are expected to
matter more for the online RNNs (daily adaptation) and TimeXer's structured
endogenous stream.

\begin{table}[h]
\centering
\small
\begin{tabular}{lrr}
\toprule
XGB feature set (200 train days) & \#feat & $\Rw$ \\
\midrule
base                          & 125 & $+0.00594$ \\
\;+ lagged responders (raw)   & 134 & $+0.00591$ \\
\;+ responder signals (engineered) & 139 & $\mathbf{+0.00624}$ \\
\bottomrule
\end{tabular}
\caption{Responder feature-engineering ablation. Raw lags are flat; the
engineered multi-scale/spread signals give a real lift.}
\end{table}

\section{Negative results and discarded approaches}
We keep an explicit record of what was tried and dropped, so the roadmap is not
re-litigated. Fuller evidence lives in \texttt{docs/WHAT\_DIDNT\_WORK.md}.
\begin{itemize}[leftmargin=1.2em,itemsep=2pt]
  \item \textbf{Signatures as ensemble members.} The log-signature / Volterra
    machinery works, and AR-based channel selection beats univariate selection,
    but no signature model moved the ensemble --- they are too correlated with
    the RNNs, which already encode the intraday path. Retained as reusable
    theory (\texttt{janestreet.theory.signatures}), not in the production blend.
  \item \textbf{Full history / curriculum.} An 800-date curriculum ($+0.00765$)
    underperformed the 280-date recency baseline ($+0.00893$); old data yields a
    negative static $R^2$ recovered only by online refit. The recency knee sits
    near ${\sim}400$ recent dates.
  \item \textbf{Symbol identity.} A 39-way \texttt{symbol\_id} one-hot was flat
    for XGB ($-0.00002$) --- the per-symbol level is already normalised away.
    Only the temporally \emph{stable} feature-profile cluster helped, modestly
    ($+0.00022$); responder co-movement clustering is unstable (early-vs-late
    ARI $-0.07$) and unusable for routing.
  \item \textbf{Raw lagged responders (static).} Flat for static XGB
    ($+0.00591$ vs.\ $+0.00594$); useful only via the ModelR aux path and the
    engineered responder signals.
  \item \textbf{Full-pool single window.} Infeasible in memory (${\sim}28$\,GB)
    and unwanted given the recency result.
  \item \textbf{Vanilla transformer and Mamba.} The vanilla transformer was
    superseded by the inverted / cross-attention variants; Mamba was built but
    never validated and has been removed. iTransformer and TimeXer remain as
    \emph{unproven} candidates.
\end{itemize}

\section{Roadmap}
In priority order: (1)~temporal bagging of the RNN models at scale;
(2)~confirm whether lagged responders help the \emph{online} models and TimeXer
even though they did not help static XGB; (3)~a TimeXer vs.\ flat-lags
head-to-head; (4)~push the training pool to the full clean range on GPU;
(5)~multi-seed single-architecture ensembles; (6)~per-architecture refit-LR
schedules. De-prioritised: larger single models, non-linear stacking, and
further signature-channel tuning, all of which showed diminishing returns.

\end{document}
